\chapter{Complexity — Time Constructible}
%%% generated by the blueprint pipeline from TCSlib.Complexity.ClassP.TimeConstructible
%%%%%%%%%%%%%%%%%%%%%%%%%%%%%%%%%%%%%%%%%%%%%%%%%%%%%%%%%%%%%%%%%%%%%%%%%%%%%%%
\section{Overview}

This module defines time-constructible functions in the sense of Arora--Barak
\S1.3: a time bound $T$ is time constructible when $T(n) \ge n$ and some machine
can produce the binary representation of $T(|x|)$ from any input $x$ within a
constant multiple of $T(|x|) + 1$ steps. The main result restores the book's
first example under this definition: the identity function is time
constructible, witnessed by a four-state machine that maintains a
little-endian binary counter on its work tape with an amortized potential
argument, then emits the counter as output.

%%%%%%%%%%%%%%%%%%%%%%%%%%%%%%%%%%%%%%%%%%%%%%%%%%%%%%%%%%%%%%%%%%%%%%%%%%%%%%%
\section{Declarations}

\begin{definition}[Time-constructible function]
\lean{Complexity.TimeConstructible}
\label{Complexity.TimeConstructible}
\leanok
\uses{Turing.FinTM, Turing.FinTM.ComputesInTime}
A function $T : \mathbb{N} \to \mathbb{N}$ is \emph{time constructible} if
$n \le T(n)$ for every $n$, and there exist a positive integer constant $c$ and
a finite Turing machine over the binary alphabet that, on every input word
$x \in \{0,1\}^{*}$, computes the binary representation of $T(|x|)$
(least-significant bit first, with no redundant leading zeros, so $0$ is
represented by the empty word) within at most $c \cdot (T(|x|) + 1)$ steps.
\end{definition}

\begin{definition}[Increment of a little-endian binary word]
\lean{Complexity.counterInc}
\label{Complexity.counterInc}
\leanok
The successor operation on binary words written least-significant bit first:
the empty word maps to the single-bit word $1$; a word beginning with $0$ has
that bit set to $1$ and the rest unchanged; a word beginning with $1$ has that
bit cleared to $0$ and the increment carries recursively into the remaining
word.
\end{definition}

\begin{definition}[Carry length of a binary word]
\lean{Complexity.counterCarry}
\label{Complexity.counterCarry}
\leanok
For a binary word written least-significant bit first, the number of leading
$1$ bits — that is, the number of cells cleared to $0$ when the word is
incremented. A word that is empty or begins with $0$ has carry length $0$.
\end{definition}

\begin{lemma}[Potential accounting for one increment]
\lean{Complexity.counterInc_potential}
\label{Complexity.counterInc_potential}
\leanok
\uses{Complexity.counterCarry, Complexity.counterInc}
\difficulty{3}
Let $b$ be a binary word written least-significant bit first, let $r$ be its
carry length (its number of leading $1$ bits), and let $w(b)$ denote the number
of $1$ bits in a word. Then the incremented word $b'$ satisfies
$w(b') + r = w(b) + 1$: each cleared bit lowers the count of ones by one, and
the single final write raises it by one.
\end{lemma}

\begin{lemma}[Increment agrees with binary successor]
\lean{Complexity.counterInc_bits}
\label{Complexity.counterInc_bits}
\leanok
\uses{Complexity.counterInc}
\difficulty{4}
For every natural number $n$, applying the word increment to the little-endian
binary representation of $n$ (with no redundant leading zeros; $0$ is the empty
word) yields exactly the binary representation of $n + 1$, including the case
where the word grows by one bit on overflow.
\end{lemma}

\begin{lemma}[Length bounds for one increment]
\lean{Complexity.counterInc_length}
\label{Complexity.counterInc_length}
\leanok
\uses{Complexity.counterCarry, Complexity.counterInc}
\difficulty{3}
For every binary word $b$ (least-significant bit first), the incremented word
has length at most $|b| + 1$, and the carry length of $b$ is at most the length
of the incremented word.
\end{lemma}

\begin{lemma}[Binary length bound]
\lean{Complexity.counter_bits_length}
\label{Complexity.counter_bits_length}
\leanok
\uses{Complexity.counterInc_bits, Complexity.counterInc_length}
\difficulty{3}
For every natural number $n$, the binary representation of $n$ (with no
redundant leading zeros) has at most $n$ bits.
\end{lemma}

\begin{definition}[Single carry transition action]
\lean{Complexity.counterBump}
\label{Complexity.counterBump}
\leanok
\uses{Turing.Action}
Given an input-head displacement $d \in \{-1, 0, +1\}$ and the symbol currently
scanned on the work tape, the action of one carry step of a one-work-tape
binary machine: if the scanned cell holds $1$, it is overwritten with $0$, the
work head moves right, and the machine continues in its carry state; otherwise
(a $0$ or a blank) the cell is overwritten with $1$, the work head moves left,
and the machine enters its rewind state. In both cases the input head moves by
$d$ and nothing is written to the output.
\end{definition}

\begin{definition}[The four-state binary-counter machine]
\lean{Complexity.counterTM}
\label{Complexity.counterTM}
\leanok
\uses{Complexity.counterBump, Turing.FinTM}
A finite Turing machine over the binary alphabet with one work tape and four
states: \emph{count}, \emph{carry}, \emph{rewind}, and \emph{emit}. In the
count state it performs a carry transition while also advancing the input head,
or switches to the emit state when the input is exhausted. In the carry state
it propagates an increment along the work tape (clearing $1$s rightward until
it writes a $1$ over the first $0$ or blank). In the rewind state it moves the
work head left until it finds the blank at cell $-1$, then steps right to cell
$0$ and re-enters the count state. In the emit state it copies the work-tape
word to the output, one bit per step moving right, and halts at the first
blank.
\end{definition}

\begin{definition}[Tape of a finite binary word]
\lean{Complexity.counterTape}
\label{Complexity.counterTape}
\leanok
The work-tape contents determined by a finite binary word $b$: for an integer
position $z$, cell $z$ holds the $z$-th bit of $b$ when $0 \le z < |b|$, and is
blank at every other position (in particular at all negative positions).
\end{definition}

\begin{definition}[Canonical counter configuration]
\lean{Complexity.counterCfg}
\label{Complexity.counterCfg}
\leanok
\uses{Complexity.counterTape, Turing.Cfg}
The configuration of the one-work-tape four-state counter machine on an input
word $x$ specified by: a state $q$ among the four counter states, an input-head
position $p$, an integer work-head position $z$, a finite binary word $b$ laid
out on the nonnegative cells of the work tape (blank elsewhere), and an output
word.
\end{definition}

\begin{lemma}[Reading past a prefix]
\lean{Complexity.counterTape_read}
\label{Complexity.counterTape_read}
\leanok
\uses{Complexity.counterTape}
\difficulty{3}
Let $p$ and $b$ be binary words, and lay the concatenation of $p$ followed by
$b$ on the nonnegative cells of a tape (blank elsewhere). Then the symbol at
cell $|p|$ is the first bit of $b$, or blank when $b$ is empty.
\end{lemma}

\begin{lemma}[Writing at the end of a prefix]
\lean{Complexity.counterTape_write}
\label{Complexity.counterTape_write}
\leanok
\uses{Complexity.counterTape, Complexity.counterTape_read}
\difficulty{4}
Let $p$ and $b$ be binary words, lay their concatenation on the nonnegative
cells of a tape, and let $\beta$ be a bit. Overwriting cell $|p|$ with $\beta$
yields exactly the tape of the word obtained by following $p$ with $\beta$ and
then the tail of $b$ (all bits of $b$ after its first). In particular, when $b$
is empty the stored word is extended by one cell; all other cells, including
the negative ones, are unchanged.
\end{lemma}

\begin{lemma}[One carry step of the counter machine]
\lean{Complexity.counter_carry_step}
\label{Complexity.counter_carry_step}
\leanok
\uses{Complexity.counterBump,
  Complexity.counterCfg,
  Complexity.counterTM,
  Complexity.counterTape,
  Complexity.counterTape_read,
  Complexity.counterTape_write,
  Turing.Action.apply,
  Turing.MultiTapeTM.step,
  Turing.moveInputPos_zero}
\difficulty{4}
Suppose the four-state binary-counter machine is in its carry state with the
concatenation of binary words $p$ and $b$ on the nonnegative work-tape cells,
work head at cell $|p|$, and empty output. If $b$ begins with $1$, one step
clears that bit to $0$ and moves the work head right, staying in the carry
state; otherwise one step writes $1$ at cell $|p|$ (replacing the first bit of
$b$, or extending the word) and moves the work head left, entering the rewind
state. The input head does not move.
\end{lemma}

\begin{lemma}[The full carry phase]
\lean{Complexity.counter_carry}
\label{Complexity.counter_carry}
\leanok
\uses{Complexity.counterCarry,
  Complexity.counterCfg,
  Complexity.counterInc,
  Complexity.counterTM,
  Complexity.counter_carry_step,
  Turing.MultiTapeTM.runFrom,
  Turing.MultiTapeTM.runFrom_succ_eq_step,
  Turing.MultiTapeTM.runFrom_zero}
\difficulty{5}
Suppose the four-state binary-counter machine is in its carry state with the
concatenation of binary words $p$ and $b$ on the nonnegative work-tape cells,
work head at cell $|p|$, and empty output, and let $r$ be the carry length of
$b$ (its number of leading $1$ bits). Then after exactly $r + 1$ steps the
machine is in the rewind state with the work tape holding $p$ followed by the
incremented word of $b$, the work head at cell $|p| + r - 1$, and the input
head and output unchanged.
\end{lemma}

\begin{lemma}[The rewind phase]
\lean{Complexity.counter_rewind}
\label{Complexity.counter_rewind}
\leanok
\uses{Complexity.counterCfg,
  Complexity.counterTM,
  Complexity.counterTape,
  Turing.Action.apply,
  Turing.Cfg.workTapeSymbols,
  Turing.MultiTapeTM.runFrom,
  Turing.MultiTapeTM.runFrom_succ_eq_step,
  Turing.MultiTapeTM.runFrom_zero,
  Turing.MultiTapeTM.step,
  Turing.moveInputPos_zero}
\difficulty{5}
Suppose the four-state binary-counter machine is in its rewind state with a
binary word $b$ on the nonnegative work-tape cells, empty output, and work head
at cell $j - 1$ for some natural number $j \le |b|$ (so possibly at the blank
cell $-1$). Then after exactly $j + 1$ steps the machine is in the count state
with the work head at cell $0$, and the stored word, input-head position, and
output all unchanged.
\end{lemma}

\begin{lemma}[The count transition consumes one input symbol]
\lean{Complexity.counter_start}
\label{Complexity.counter_start}
\leanok
\uses{Complexity.counterBump,
  Complexity.counterCfg,
  Complexity.counterTM,
  Complexity.counterTape,
  Turing.Action.apply,
  Turing.Cfg.inputSymbol,
  Turing.inputSymbolInner,
  Turing.moveInputPos,
  Turing.moveInputPos_pos_of_ne_right,
  Turing.moveInputPos_zero}
\difficulty{4}
Let $x$ be an input word and $i < |x|$. Taking one step of the four-state
binary-counter machine from the count state — with the input head scanning the
$i$-th input symbol, the work head at cell $0$ over a stored binary word $b$,
and empty output — produces the same configuration as taking one step from the
carry state with the input head advanced by one position and the same work tape
and output.
\end{lemma}

\begin{lemma}[Cost of one complete increment]
\lean{Complexity.counter_increment}
\label{Complexity.counter_increment}
\leanok
\uses{Complexity.counterCarry,
  Complexity.counterCfg,
  Complexity.counterInc,
  Complexity.counterInc_length,
  Complexity.counterTM,
  Complexity.counter_carry,
  Complexity.counter_rewind,
  Complexity.counter_start,
  Turing.MultiTapeTM.runFrom,
  Turing.MultiTapeTM.runFrom_add,
  Turing.MultiTapeTM.runFrom_succ_eq_step}
\difficulty{4}
Let $x$ be an input word and $i < |x|$. Starting the four-state binary-counter
machine in the count state with the input head scanning the $i$-th input
symbol, a binary word $b$ stored on the work tape with the work head at cell
$0$, and empty output: after exactly $2r + 2$ steps, where $r$ is the carry
length of $b$, the machine is again in the count state with the input head
advanced by one position, the work tape holding the incremented word of $b$,
the work head back at cell $0$, and the output still empty.
\end{lemma}

\begin{lemma}[The amortized counting invariant]
\lean{Complexity.counter_count}
\label{Complexity.counter_count}
\leanok
\uses{Complexity.counterCarry,
  Complexity.counterCfg,
  Complexity.counterInc_bits,
  Complexity.counterInc_potential,
  Complexity.counterTM,
  Complexity.counterTape,
  Complexity.counter_increment,
  Turing.MultiTapeTM.initCfg,
  Turing.MultiTapeTM.runFrom,
  Turing.MultiTapeTM.runFrom_add}
\difficulty{4}
Let $x$ be a binary input word and $i \le |x|$, and let $w(i)$ denote the
number of $1$ bits in the binary representation of $i$. Then there is a step
count $t$ with $t + 2\,w(i) \le 4i$ such that, after $t$ steps from its initial
configuration on $x$, the four-state binary-counter machine is in the count
state having consumed exactly the first $i$ input symbols, with the work tape
holding the binary representation of $i$ (least-significant bit first), the
work head at cell $0$, and empty output.
\end{lemma}

\begin{lemma}[The emit phase copies the word]
\lean{Complexity.counter_emit_run}
\label{Complexity.counter_emit_run}
\leanok
\uses{Complexity.counterCfg,
  Complexity.counterTM,
  Complexity.counterTape,
  Turing.Action.apply,
  Turing.Cfg.workTapeSymbols,
  Turing.MultiTapeTM.runFrom,
  Turing.MultiTapeTM.runFrom_succ_eq_step',
  Turing.MultiTapeTM.step,
  Turing.moveInputPos_zero}
\difficulty{4}
Suppose the four-state binary-counter machine is in its emit state with a
binary word $b$ on the nonnegative work-tape cells, work head at cell $0$, and
empty output. Then for every $i \le |b|$, after exactly $i$ steps it is still
in the emit state with the work head at cell $i$, the tape contents and
input-head position unchanged, and the output equal to the first $i$ bits of
$b$.
\end{lemma}

\begin{lemma}[Emission halts with the stored word]
\lean{Complexity.counter_emit}
\label{Complexity.counter_emit}
\leanok
\uses{Complexity.counterCfg,
  Complexity.counterTM,
  Complexity.counterTape,
  Complexity.counter_emit_run,
  Turing.Action.apply,
  Turing.Cfg.workTapeSymbols,
  Turing.MultiTapeTM.runFrom,
  Turing.MultiTapeTM.runFrom_succ_eq_step',
  Turing.MultiTapeTM.step}
\difficulty{3}
Suppose the four-state binary-counter machine is in its emit state with a
binary word $b$ on the nonnegative work-tape cells, work head at cell $0$, and
empty output. Then after exactly $|b| + 1$ steps the machine has halted and its
output equals $b$.
\end{lemma}

\begin{theorem}[The identity function is time constructible]
\lean{Complexity.timeConstructible_id}
\label{Complexity.timeConstructible_id}
\leanok
\uses{Complexity.TimeConstructible,
  Complexity.counterCfg,
  Complexity.counterTM,
  Complexity.counter_bits_length,
  Complexity.counter_count,
  Complexity.counter_emit,
  Turing.Action.apply,
  Turing.Cfg.inputSymbol,
  Turing.FinTM.ComputesInTime,
  Turing.FinTM.ComputesInTime.mono,
  Turing.MultiTapeTM.initCfg,
  Turing.MultiTapeTM.runFrom,
  Turing.MultiTapeTM.runFrom_add,
  Turing.MultiTapeTM.runFrom_succ_eq_step',
  Turing.MultiTapeTM.step}
\difficulty{5}
The identity function $n \mapsto n$ on the natural numbers is time
constructible: trivially $n \le n$ for all $n$, and there are a positive
integer constant $c$ and a finite Turing machine over the binary alphabet that,
on every input word $x \in \{0,1\}^{*}$, outputs the binary representation of
$|x|$ (least-significant bit first, empty for $|x| = 0$) within at most
$c \cdot (|x| + 1)$ steps.
\end{theorem}
